\documentclass[12pt]{article}
\usepackage{amsmath, amssymb, amsthm, geometry, mathtools}
\usepackage{hyperref}
\geometry{a4paper, margin=1in}

\title{Nonlinearity and Linearity: A Unified Bundle Perspective on Finite and Infinite-Dimensional Dynamics}
\author{Shimon Panfil}
\date{\today}

\theoremstyle{definition}
\newtheorem{definition}{Definition}[section]
\newtheorem{example}{Example}[section]

\theoremstyle{plain}
\newtheorem{theorem}{Theorem}[section]
\newtheorem{proposition}{Proposition}[section]
\newtheorem{corollary}{Corollary}[section]

\begin{document}

\maketitle

\begin{abstract}
We present a unified framework relating nonlinear finite-dimensional dynamics to linear infinite-dimensional evolution, and conversely, linear infinite-dimensional systems to effective finite-dimensional nonlinear dynamics. Using vector bundles and sections as a geometric language, we unify Koopman linearization, Carleman embedding, and Naimark dilation, as well as projection-induced nonlinearities such as mean-field reductions and moment truncations. This framework clarifies the structural origins of nonlinearity in both classical and quantum systems.
\end{abstract}

\section{Introduction}

Nonlinear dynamics in finite-dimensional systems are ubiquitous in physics and mathematics. Remarkably, these dynamics can often be \emph{linearized} by embedding into an infinite-dimensional space of functions, sections, or extended Hilbert spaces. Conversely, infinite-dimensional linear systems often produce effective nonlinear dynamics when restricted to finite-dimensional manifolds or observables. We unify these perspectives using bundles and sections.

\section{Bundles, Sections, and Observables}

\begin{definition}[Vector bundle]
A \emph{vector bundle} $\pi: E \to M$ over a smooth manifold $M$ consists of a smooth manifold $E$ such that each fiber $E_x := \pi^{-1}(x)$ is a vector space and locally $E|_U \simeq U \times \mathbb{C}^k$.
\end{definition}

\begin{definition}[Section]
A \emph{section} of $E$ is a smooth map $s: M \to E$ satisfying $\pi \circ s = \mathrm{id}_M$. The space of smooth sections is denoted $\Gamma(E)$.
\end{definition}

Classical observables are sections of the trivial line bundle $M \times \mathbb{C}$, and quantum states may be viewed as sections of Hilbert bundles $\mathcal{H}_M \to M$.

\section{Nonlinear Finite-Dimensional $\to$ Linear Infinite-Dimensional}

Let $\Phi_t: M \to M$ be a smooth flow generated by a vector field $F \in \Gamma(TM)$.

\begin{definition}[Pullback on sections]
The pullback operator $U_t: \Gamma(E) \to \Gamma(E)$ is defined as
\[
(U_t s)(x) := s(\Phi_t(x)).
\]
\end{definition}

\begin{theorem}[Bundle Koopman Linearization]
$U_t$ is linear, forms a semigroup $U_{t+s} = U_t \circ U_s$, and its generator is the Lie derivative $L = \mathcal{L}_F$ along $F$.
\end{theorem}

\begin{proof}
Linearity and semigroup properties follow from the definition. Differentiating along the flow yields $L s = \mathcal{L}_F s$.
\end{proof}

\begin{example}[Carleman Linearization]
A polynomial ODE $\dot{x} = P(x)$ with $x \in \mathbb{R}^n$ can be lifted to a linear operator on the symmetric tensor bundle $S^\bullet(T^*\mathbb{R}^n)$. Projection to degree-one monomials recovers the original nonlinear dynamics.
\end{example}

\begin{example}[Naimark Dilation]
Let $E: \Sigma \to \mathcal{B}(\mathcal{H})$ be a POVM. Naimark's theorem provides a Hilbert space $\mathcal{K} \supset \mathcal{H}$, a PVM $\tilde{E}$, and isometry $V: \mathcal{H} \to \mathcal{K}$ such that
\[
E(A) = V^\dagger \tilde{E}(A) V, \quad \forall A \in \Sigma.
\]
Nonlinear POVM collapse becomes linear evolution on the extended bundle.
\end{example}

\section{Linear Infinite-Dimensional $\to$ Nonlinear Finite-Dimensional}

Let $V$ be an infinite-dimensional vector space with linear evolution $L: V \to V$. Let $M \subset V$ be a finite-dimensional embedded submanifold with smooth projection $\pi: V \to M$.

\begin{theorem}[Projection-Induced Nonlinearity]
The induced dynamics on $M$,
\[
\dot{x} = \pi(L x),
\]
is generally nonlinear.
\end{theorem}

\begin{proof}
The projection $\pi$ is nonlinear in general; applying it to the linear flow $L$ yields nonlinear evolution in the coordinates of $M$.
\end{proof}

\begin{example}[Moment Truncation]
From the Liouville equation $\partial_t \rho + \nabla \cdot (F \rho) = 0$, projecting onto the first moments $\langle x \rangle(t)$ gives a nonlinear ODE for $\langle x \rangle$.
\end{example}

\begin{example}[Mean-Field Quantum Reduction]
Linear many-body Schrödinger evolution on $\mathcal{H}^{\otimes N}$ reduces to a nonlinear mean-field equation (e.g., Gross-Pitaevskii) on the product manifold $\psi^{\otimes N}$.
\end{example}

\section{Unified Bundle Perspective}

Both directions can be expressed geometrically:

\begin{itemize}
\item Finite nonlinear dynamics $\Phi_t$ $\implies$ linear pullback $U_t$ on $\Gamma(E)$.
\item Infinite linear evolution $L$ on $V$ $\implies$ nonlinear projection $\pi(L)$ on $M \subset V$.
\end{itemize}

Symbolically:
\[
\text{Nonlinear finite} \leftrightarrow \text{Linear infinite}, \quad
\text{Linear infinite} \leftrightarrow \text{Nonlinear finite}.
\]

Projection or restriction is the geometric source of nonlinearity, while embedding or pullback restores linearity.

\section{Discussion}

This framework unifies:

\begin{enumerate}
\item Koopman linearization (classical flows)
\item Carleman embedding (polynomial ODEs)
\item Naimark dilation (quantum POVMs)
\item Projection-induced nonlinearities (mean-field, moment truncation)
\end{enumerate}

Nonlinearity is seen as an artifact of projecting from a larger linear bundle or Hilbert space. Infinite-dimensional linear spaces are natural universes for dynamics; finite-dimensional nonlinear behavior emerges from restriction.

\section{Conclusion}

Viewing dynamics through bundles and sections provides a unified language connecting classical and quantum linearizations, as well as projection-induced nonlinearities. Both directions—linearization of finite-dimensional nonlinear systems and reduction of infinite-dimensional linear systems—share the same geometric principle: **nonlinearity arises from truncation or projection, linearity resides in the extended space**.

\begin{thebibliography}{9}

\bibitem{Koopman1931}
B. O. Koopman, \emph{Hamiltonian Systems and Transformation in Hilbert Space}, Proc. Natl. Acad. Sci. USA, 17 (1931), 315--318.

\bibitem{Naimark1940}
M. A. Naimark, \emph{Spectral Functions of a Symmetric Operator}, Izv. Akad. Nauk SSSR Ser. Mat. 4 (1940), 277–318.

\end{thebibliography}

\end{document}
